\section*{Chapitre \ref{ch:g11:nucleus-radioactivity} --- Le noyau et la radioactivité}

\begin{solution}{exo:g11:nucleus-radioactivity:1}
${}^{16}_{8}\mathrm{O}$~: $8$ p, $8$ n~; ${}^{27}_{13}\mathrm{Al}$~: $13$ p,
$14$ n~; ${}^{60}_{27}\mathrm{Co}$~: $27$ p, $33$ n~;
${}^{235}_{92}\mathrm{U}$~: $92$ p, $143$ n~; ${}^{3}_{1}\mathrm{H}$~: $1$ p,
$2$ n. ($N = A - Z$ chaque fois.)
\end{solution}

\begin{solution}{exo:g11:nucleus-radioactivity:2}
Les trois carbones ($Z = 6$) sont des
\omterm{def:g11:nucleus-radioactivity:notation}{isotopes} les uns des autres.
${}^{14}_{6}\mathrm{C}$ et ${}^{14}_{7}\mathrm{N}$ partagent $A = 14$
mais non $Z$~: éléments différents (nuages électroniques différents),
donc pas des \omterm{def:g11:nucleus-radioactivity:notation}{isotopes}.
\end{solution}

\begin{solution}{exo:g11:nucleus-radioactivity:3}
$A = 210 - 4 = 206$, $Z = 84 - 2 = 82$~:
${}^{210}_{84}\mathrm{Po} \to {}^{206}_{82}\mathrm{Pb} +
{}^{4}_{2}\mathrm{He}$ --- le \omterm{def:g11:fundamental-interactions:ladder}{noyau} fils est le plomb-206.
\end{solution}

\begin{solution}{exo:g11:nucleus-radioactivity:4}
${}^{60}_{27}\mathrm{Co} \to {}^{60}_{28}\mathrm{Ni} +
{}^{\;0}_{-1}\mathrm{e}$ et
${}^{131}_{53}\mathrm{I} \to {}^{131}_{54}\mathrm{Xe} +
{}^{\;0}_{-1}\mathrm{e}$~: $A$ inchangé, $Z$ augmente d'un.
\end{solution}

\begin{solution}{exo:g11:nucleus-radioactivity:5}
$16$~jours $= 2\,T_{1/2}$~: $800/4 = \qty{200}{MBq}$. $32$~jours
$= 4\,T_{1/2}$~: $800/16 = \qty{50}{MBq}$.
\end{solution}

\begin{solution}{exo:g11:nucleus-radioactivity:6}
(a)~$\alpha$ (lourde, donc à peine déviée)~; (b)~$\beta^-$~; (c)~$\gamma$
(neutre, pénétrante)~; (d)~$\beta^+$
(\omterm{rem:g11:nucleus-radioactivity:betaplus}{positon}~: aussi léger que
(b), charge opposée).
\end{solution}

\begin{solution}{exo:g11:nucleus-radioactivity:7}
${}^{222}_{86}\mathrm{Rn} \to {}^{218}_{84}\mathrm{Po} +
{}^{4}_{2}\mathrm{He}$~; \
${}^{218}_{84}\mathrm{Po} \to {}^{214}_{82}\mathrm{Pb} +
{}^{4}_{2}\mathrm{He}$~; \
${}^{214}_{82}\mathrm{Pb} \to {}^{214}_{83}\mathrm{Bi} +
{}^{\;0}_{-1}\mathrm{e}$.
\end{solution}

\begin{solution}{exo:g11:nucleus-radioactivity:8}
Le fluor-18 a $9$ neutrons là où le fluor-19 stable en a $10$~: pauvre
en neutrons, sous la bande, donc $\beta^+$~:
${}^{18}_{9}\mathrm{F} \to {}^{18}_{8}\mathrm{O} + {}^{0}_{+1}\mathrm{e}$.
\end{solution}

\begin{solution}{exo:g11:nucleus-radioactivity:9}
$24$~heures $= 4\,T_{1/2}$~: $500/16 \approx \qty{31}{MBq}$. Comme
$2^9 = 512$, neuf \omterm{def:g11:nucleus-radioactivity:halflife}{demi-vies} amènent $500/512 < \qty{1}{MBq}$~: après
$9 \times 6 = \qty{54}{h}$.
\end{solution}

\begin{solution}{exo:g11:nucleus-radioactivity:10}
$\frac18 = \frac1{2^3}$~: trois \omterm{def:g11:nucleus-radioactivity:halflife}{demi-vies}, $3 \times \num{5700} =
\num{17100}$~ans. Un os de dinosaure a environ $\num{17500}$ \omterm{def:g11:nucleus-radioactivity:halflife}{demi-vies}
d'âge~: une fraction $1/2^{17500}$ reste --- pas un seul
\omterm{def:g11:fundamental-interactions:ladder}{atome} de carbone-14 (un gramme de
carbone n'en détient qu'environ $\num{5e22}$
\omterm{def:g11:fundamental-interactions:ladder}{atomes}).
\end{solution}

\begin{solution}{exo:g11:nucleus-radioactivity:11}
${}^{14}_{6}\mathrm{C}$~: $8$ neutrons contre $6$--$7$ dans le carbone
stable, au-dessus de la bande, $\beta^-$. ${}^{11}_{6}\mathrm{C}$~:
$5$ neutrons, pauvre en neutrons, $\beta^+$.
${}^{238}_{92}\mathrm{U}$~: $Z = 92 > 83$, trop lourd, $\alpha$.
\end{solution}

\begin{solution}{exo:g11:nucleus-radioactivity:12}
Seule $\alpha$ change $A$~: $238 - 206 = 32 = 4x$, donc $x = 8$
\omterm{def:g11:nucleus-radioactivity:abg}{alphas}. Elles retirent $16$ de $Z$~:
$92 - 16 + y = 82$ donne $y = 6$ désintégrations bêta-moins.
\end{solution}

\begin{solution}{exo:g11:nucleus-radioactivity:13}
$6.4/0.1 = 64 = 2^{6}$~: six \omterm{def:g11:nucleus-radioactivity:halflife}{demi-vies}, $6 \times 5.3 = 31.8 \approx 32$
années en stockage blindé.
\end{solution}

\begin{solution}{exo:g11:nucleus-radioactivity:14}
Par jour~: $8000 \times 86400 \approx \num{6.9e8}$ désintégrations. Sur
$80$~ans~: $\num{6.9e8} \times 365 \times 80 \approx \num{2e13}$. La vie
a toujours tourné sur ce fond --- un
\omterm{def:g11:nucleus-radioactivity:activity}{becquerel} est un
\omterm{def:g11:fundamental-interactions:ladder}{atome} par seconde sur les
$\sim 10^{27}$ du corps~; l'\omterm{def:g10:orders-of-magnitude:unit}{unité} est
minuscule, et un décompte de
\omterm{def:g11:nucleus-radioactivity:activity}{becquerels} qui sonne effrayant peut
valoir quelques bananes.
\end{solution}

\begin{solution}{exo:g11:nucleus-radioactivity:15}
$\num{25000} \times 86400 \approx \num{2.2e9}$ désintégrations par jour.
Dix ans sont $10/432 \approx 2\,\%$ d'une
\omterm{def:g11:nucleus-radioactivity:halflife}{demi-vie}~: bien moins de la moitié a
disparu, donc l'électronique et la poussière, non la source,
imposent le remplacement. À l'extérieur, les
$\alpha$ meurent dans l'air et le boîtier de la chambre~; la poussière
inhalée gare l'émetteur dans les poumons, où chaque $\alpha$ déverse
son \omterm{def:g10:energy-conservation:energy}{énergie} dans le tissu vivant --- une
grande \omterm{rem:g11:nucleus-radioactivity:dose}{dose}
(\omterm{rem:g11:nucleus-radioactivity:dose}{sieverts}) pour une
\omterm{def:g11:nucleus-radioactivity:activity}{activité} modeste.
\end{solution}

\begin{solution}{pb:g11:nucleus-radioactivity:1}
\textbf{1.} ${}^{12}_{6}\mathrm{C}$~: $6$ p, $6$ n~; ${}^{14}_{6}\mathrm{C}$~:
$6$ p, $8$ n. La chimie ne voit que le nuage électronique, fixé par
$Z = 6$~: comportement identique.

\textbf{2.} Manger et respirer renouvellent constamment le carbone d'un
corps vivant à la proportion atmosphérique~; à la mort l'apport
s'arrête et la désintégration court sans frein.

\textbf{3.} ${}^{14}_{6}\mathrm{C} \to {}^{14}_{7}\mathrm{N} +
{}^{\;0}_{-1}\mathrm{e}$ ($\beta^-$). Ces $\beta$ sont mous~: des
\omterm{def:g10:orders-of-magnitude:unit}{mètres} de lin enveloppé les arrêtent, donc
peu s'échappe de la momie --- d'où la mesure sur un petit échantillon
de carbone prélevé sur l'objet lui-même.

\textbf{4.} $6.8/13.6 = \frac12$~: une
\omterm{def:g11:nucleus-radioactivity:halflife}{demi-vie}. La momie a environ
\num{5700}~ans.

\textbf{5.} $13.2/13.6 \approx 0.97$~: le lin est essentiellement
moderne --- au plus quelques siècles. Un faux.

\textbf{6.} $13.6/2^{10} = 13.6/1024 \approx 0.013$ désintégrations par
minute --- un coup toutes les $75$~minutes par gramme, noyé dans le fond
naturel. Dix \omterm{def:g11:nucleus-radioactivity:halflife}{demi-vies}, environ \num{57000}~ans, sont l'horizon
pratique.

\textbf{7.} $92$ protons, $146$ neutrons. La répulsion de Coulomb agit
entre toutes les paires de protons à travers le
\omterm{def:g11:fundamental-interactions:ladder}{noyau} tandis que la colle forte ne
lie que les voisins~: les noyaux lourds ne survivent qu'avec des
neutrons en surplus, colle sans répulsion.

\textbf{8.} ${}^{238}_{92}\mathrm{U} \to {}^{234}_{90}\mathrm{Th} +
{}^{4}_{2}\mathrm{He}$.

\textbf{9.} ${}^{234}_{90}\mathrm{Th} \to {}^{234}_{91}\mathrm{Pa} +
{}^{\;0}_{-1}\mathrm{e}$, puis ${}^{234}_{91}\mathrm{Pa} \to
{}^{234}_{92}\mathrm{U} + {}^{\;0}_{-1}\mathrm{e}$~: l'uranium réapparaît,
comme uranium-234.

\textbf{10.} $A$~: $238 - 206 = 32$, donc $8$ étapes $\alpha$~; $Z$~:
$92 - 16 + y = 82$, donc $6$ étapes $\beta^-$.

\textbf{11.} $\num{4.5e9}$~ans est une
\omterm{def:g11:nucleus-radioactivity:halflife}{demi-vie}~: $\frac12$ reste. Cinq
milliards d'années de plus $\approx$ une
\omterm{def:g11:nucleus-radioactivity:halflife}{demi-vie} de plus~: environ
$\frac14$.

\textbf{12.} Elle pilote les volcans, la tectonique des plaques
(continents dérivants, séismes) et la chaleur géothermique~: la
géologie de surface d'une planète chauffée de l'intérieur.

\textbf{13.} ${}^{241}_{95}\mathrm{Am} \to {}^{237}_{93}\mathrm{Np} +
{}^{4}_{2}\mathrm{He}$.

\textbf{14.} $\num{3.3e4}$ désintégrations par seconde~;
$\num{33000} \times 86400 \approx \num{2.9e9}$ par jour.

\textbf{15.} Chaque $\alpha$ ionise l'air qu'elle traverse~; les ions
portent un petit courant entre les électrodes de la chambre. Les
particules de fumée capturent les ions, le courant chute, l'alarme
sonne.

\textbf{16.} Les $\alpha$ dépensent toute leur
\omterm{def:g10:energy-conservation:energy}{énergie} dans la chambre et rien
n'échappe du boîtier~: ionisation maximale, zéro fuite. Une source
$\gamma$ d'\omterm{def:g11:nucleus-radioactivity:activity}{activité} égale ioniserait à
peine l'air de la chambre et irradierait toute la pièce à la place.

\textbf{17.} Dix ans sont $10/432 \approx 2\,\%$ d'une
\omterm{def:g11:nucleus-radioactivity:halflife}{demi-vie}~: l'
\omterm{def:g11:nucleus-radioactivity:activity}{activité} est essentiellement
inchangée. Les maillons faibles sont la poussière, les insectes et
l'électronique --- le
\omterm{def:g11:fundamental-interactions:ladder}{noyau} survit au gadget.

\textbf{18.} $\num{9e9}$~ans $= 2$ \omterm{def:g11:nucleus-radioactivity:halflife}{demi-vies}~: $\frac14$ survit.
L'uranium de la Terre est encore abondant, donc il a été forgé il y a
au plus quelques \omterm{def:g11:nucleus-radioactivity:halflife}{demi-vies} --- la planète s'est condensée à partir de
cendres d'étoiles fraîches, des milliards d'années après les premières
étoiles.

\textbf{19.} Géologiquement morte~: un intérieur froid, pas de
volcanisme ni de tectonique des plaques pour recycler l'air et la
roche, et pas de \omterm{def:g11:fundamental-interactions:ladder}{noyau} en brassage --- d'où pas de bouclier magnétique
contre le vent solaire.

\textbf{20.} Chaque \omterm{def:g11:fundamental-interactions:ladder}{atome} de vous
au-delà de l'hélium a été forgé dans une étoile mourante, et les
restes instables de cette forge chauffent encore le sol sous vos
pieds~: poussière d'étoiles, maintenue en vie par sa propre
\omterm{def:g11:nucleus-radioactivity:halflife}{demi-vie}.
\end{solution}
